\begin{abstract}
While current software agents powered by large language models (LLMs) and reinforcement learning (RL) can boost programmer productivity, their reliance on human-curated training data and environments creates a fundamental barrier to superintelligence.
In this paper, we present \mbox{\techfull} (\tech), a first step toward training superintelligent software agents under minimal data assumptions.
\tech requires only access to sandboxed repositories with source code and dependencies, no need for human-labeled issues or test commands.
Grounded in real-world codebases, a single LLM agent is trained via RL in a self-play setting to inject and repair increasingly complex bugs.
The bugs are formally specified by test suite improvements proposed by the agent rather than natural language issue descriptions.
On the \sbv and \sbp benchmarks, \tech achieves clear self-improvement (+10.4 and +7.8 points) and consistently outperforms the human-data baseline throughout training, generalizing to natural language bug descriptions not seen in training.
Overall, our results point toward a paradigm where agents autonomously gather extensive learning experiences from real software repositories, ultimately enabling superintelligent systems that exceed human capabilities in understanding, modifying, and creating software from scratch.


\end{abstract}
